%% ---------------------------------------------------------------------
\section{Tipo del sistema y constantes de error}
\label{sec:er-tipo}
%% ---------------------------------------------------------------------

Con retroalimentación unitaria (\cref{fig:er-unitario}), $\Tlc=G/(1+G)$ y $E(s)=\dfrac{1}{1+G(s)}R(s)$. Si el lazo cerrado es estable, el teorema del valor final da
\begin{equation}
  \boxed{\;\ess=\lim_{t\to\infty}e(t)=\lim_{s\to0}\,s\,\frac{R(s)}{1+G(s)}\;}
  \label{eq:er-ess}
\end{equation}
Aquí $G(s)$ es el \textbf{lazo abierto}: la función que, con retroalimentación unitaria, va del error a la salida (\cref{def:rf-malla}; si la retroalimentación no es unitaria, ver la \cref{sec:er-no-unitaria}).

\begin{figure}[ht]
\centering
\includegraphics[width=0.55\linewidth]{cap06/lazo_unitario}
\caption{Lazo con retroalimentación unitaria.}
\label{fig:er-unitario}
\end{figure}

Escribamos el lazo abierto en forma de Bode, separando los polos en el origen:
\begin{equation}
  G(s)=\frac{K\,(T_1s+1)(T_2s+1)\cdots(T_ms+1)}{s^N(\tau_1s+1)(\tau_2s+1)\cdots(\tau_{n-N}s+1)} .
  \label{eq:er-bode}
\end{equation}

\begin{definicion}[label=def:er-tipo]{Tipo del sistema}
El \textbf{tipo} es el número $N$ de polos en el origen del lazo abierto, es decir, el número de integradores del lazo. $K$ es su ganancia de Bode: el valor de $s^NG(s)$ en $s=0$.
\end{definicion}

Cada señal de prueba define una \textbf{constante estática de error}; cuanto mayor, menor el error:
\begin{itemize}
  \item \textbf{Escalón}, $R=1/s$: $\ess=\dfrac{1}{1+\lim_{s\to0}G(s)}=\dfrac{1}{1+\kp}$, con la \clave{constante de posición} $\kp=\lim_{s\to0}G(s)$.
  \item \textbf{Rampa}, $R=1/s^2$: $\ess=\lim_{s\to0}\dfrac{1}{s+sG(s)}=\dfrac{1}{\kv}$, con la \clave{constante de velocidad} $\kv=\lim_{s\to0}sG(s)$.
  \item \textbf{Parábola}, $R=1/s^3$: $\ess=\dfrac{1}{\ka}$, con la \clave{constante de aceleración} $\ka=\lim_{s\to0}s^2G(s)$.
\end{itemize}
Con \eqref{eq:er-bode}, $s^iG(s)\to K/s^{N-i}$: la constante vale $K$ si $N=i$, infinito si $N>i$ y cero si $N<i$. Eso da la tabla que resume el capítulo.

\begin{table}[H]
\centering
\caption{Error en estado estacionario según el tipo del lazo abierto (lazo cerrado estable, retroalimentación unitaria). Para una entrada de amplitud $A$ (escalón $A$, rampa $At$, parábola $At^2/2$), el error se multiplica por $A$.}
\label{tab:er-tabla}
\renewcommand{\arraystretch}{1.5}
\begin{tabular}{@{}lcccccc@{}}
\toprule
 & $\kp$ & $\kv$ & $\ka$ & escalón & rampa & parábola\\
\midrule
Tipo 0 & $K$ & $0$ & $0$ & $\dfrac{1}{1+\kp}$ & $\infty$ & $\infty$\\
Tipo 1 & $\infty$ & $K$ & $0$ & $0$ & $\dfrac{1}{\kv}$ & $\infty$\\
Tipo 2 & $\infty$ & $\infty$ & $K$ & $0$ & $0$ & $\dfrac{1}{\ka}$\\
\bottomrule
\end{tabular}
\end{table}

\begin{ideaclave}
\textbf{Cada integrador vale una derivada de la referencia.}\enspace Un lazo de tipo $N$ sigue sin error una referencia polinomial de grado menor que $N$, con error finito una de grado $N$, y pierde una de grado mayor. El error se lee en la tabla con solo dos datos: el tipo y una ganancia.
\end{ideaclave}

\begin{figure}[ht]
\centering
\includegraphics[width=\linewidth]{cap06/f_tipos}
\caption{Lazos abiertos de tipo 0, 1 y 2 ($\frac{4}{(s+1)(s+2)}$, $\frac{4}{s(s+2)}$ y $\frac{4(s+0.5)}{s^2(s+2)}$) ante escalón y rampa unitarios. Errores: $1/3$ y $\infty$; $0$ y $0.5$; $0$ y $0$.}
\label{fig:er-tipos}
\end{figure}

\begin{ejemplo}[label=ej:er-motor]{Un motor con retroalimentación de velocidad}
En el servo de la \cref{fig:er-motor}~\citep{DoBi08}, halle $\ess$ ante una rampa unitaria en función de $K$, $K_b$ y $K_m$, y escoja $K$ para $\ess=1$ con $K_m=10$ y $K_b=0.05$.
\solucion El lazo interno da $\Omega/(KE)=K_m/(s+0.01+K_bK_m)$, y con el integrador,
\[
  G(s)=\frac{KK_m}{s(s+0.01+K_bK_m)} :
\]
tipo 1, con $\kv=\dfrac{KK_m}{0.01+K_bK_m}$ y $\ess=\dfrac{0.01+K_bK_m}{KK_m}$. Con los datos, $\ess=0.51/(10K)=1$ da $K=0.051$. (El lazo cerrado es de segundo orden con coeficientes positivos: estable para todo $K>0$.)
\end{ejemplo}

\begin{figure}[ht]
\centering
\includegraphics[width=0.85\linewidth]{cap06/lazo_motor_kb}
\caption{Ejemplo~\ref{ej:er-motor}: servo de posición con retroalimentación interna de velocidad.}
\label{fig:er-motor}
\end{figure}

\codigorepo{cap06_error/C6_Ej1_TipoError.m}{tipo}{tipo, constantes y estabilidad de cuatro lazos con la función \cmd{constantesError}}
\codigorepo{cap06_error/C6_Ej1_TipoError.m}{simulacion}{el error de rampa simulado}

\begin{pausa}
Un lazo abierto $G=\dfrac{5(s+1)}{s^3(s+2)}$ tiene $\kp=\kv=\ka=\infty$. ¿Sigue sin error una parábola?
\respuesta{No se puede saber con la tabla: el polinomio característico es $s^4+2s^3+5s+5$, al que le falta el término en $s^2$, así que el lazo cerrado es inestable. La salida no sigue nada.}
\end{pausa}
